\section{Sensor Validation}
\label{sec:ch4-sensor-validation}

The sensor-validation results determine whether the calibrated measurements
were reliable enough for the protection decisions described in
Chapter~\ref{chap:methodology}. In this section, the voltage, current, and
device-temperature channels are interpreted in terms of decision reliability rather
than as isolated measurement exercises.

\subsection{Voltage Sensor Validation}
\label{sec:ch4-voltage-validation}

The voltage sensor was compared with a Fluke 115 True RMS multimeter at ten
test points.
\begin{table}[H]
  \centering
  \caption{Validation results for the voltage sensor}
  \label{tab:ch4-voltage-validation}
  \begin{tabular}{ccccc}
    \toprule
    \textbf{Test} & \textbf{Device (V)} & \textbf{Fluke (V)} & \textbf{Abs. error (V)} & \textbf{Squared error (V$^2$)} \\
    \midrule
    1 & 181.0 & 184.7 & 3.7 & 13.69 \\
    2 & 196.1 & 192.8 & 3.3 & 10.89 \\
    3 & 198.4 & 201.9 & 3.5 & 12.25 \\
    4 & 210.3 & 213.6 & 3.3 & 10.89 \\
    5 & 223.8 & 219.7 & 4.1 & 16.81 \\
    6 & 221.0 & 225.4 & 4.4 & 19.36 \\
    7 & 234.6 & 230.9 & 3.7 & 13.69 \\
    8 & 232.7 & 236.5 & 3.8 & 14.44 \\
    9 & 247.5 & 243.1 & 4.4 & 19.36 \\
    10 & 244.3 & 249.2 & 4.9 & 24.01 \\
    \bottomrule
  \end{tabular}
\end{table}

Table~\ref{tab:ch4-voltage-validation} shows the validation results for the voltage sensor. 
The corrected prototype readings produced an MAE of 3.910V, MSE
of 15.539V$^2$, and RMSE of 3.94V. These errors are small relative to the
firmware voltage bands in Table~\ref{tab:ch3-voltage-thresholds}; therefore,
the calibrated voltage channel is adequate for distinguishing normal operation
from the undervoltage and overvoltage regions used in this prototype.
The result supports voltage
state discrimination, but it does not by itself certify the device for
utility-grade metering.

\subsection{Current Sensor Validation}
\label{sec:ch4-current-validation}

The current sensor was compared with an Ingco clamp meter.
\begin{table}[H]
  \centering
  \caption{Validation results for the current sensor}
  \label{tab:ch4-current-validation}
  \begin{tabular}{ccccc}
    \toprule
    \textbf{Test} & \textbf{Device (A)} & \textbf{Clamp meter (A)} & \textbf{Abs. error (A)} & \textbf{Squared error (A$^2$)} \\
    \midrule
    1 & 0.66 & 0.68 & 0.02 & 0.0004 \\
    2 & 1.98 & 1.94 & 0.04 & 0.0016 \\
    3 & 5.00 & 5.11 & 0.11 & 0.0121 \\
    4 & 7.80 & 7.63 & 0.17 & 0.0289 \\
    5 & 8.97 & 9.14 & 0.17 & 0.0289 \\
    6 & 10.76 & 10.58 & 0.18 & 0.0324 \\
    7 & 11.84 & 12.06 & 0.22 & 0.0484 \\
    8 & 14.00 & 13.74 & 0.26 & 0.0676 \\
    9 & 14.34 & 14.67 & 0.33 & 0.1089 \\
    10 & 16.14 & 15.86 & 0.28 & 0.0784 \\
    \bottomrule
  \end{tabular}
\end{table}

Table~\ref{tab:ch4-current-validation} shows that error increased slightly at
higher currents but remained much smaller than the distance between the warning
and sustained-trip thresholds.
The current channel
produced an MAE of 0.178A, MSE of 0.04076A$^2$, and RMSE of 0.202A. This
level of error is acceptable for the implemented warning and sustained-trip
thresholds because the warning threshold is 10.0A and the sustained-overload
threshold is 14.5A. The same current channel is also used by the socket
thermal model and the TinyML feature extractor, so a low RMS-current error
improves both threshold and feature reliability.
This supports the overload results discussed in
Section~\ref{sec:ch4-overload-protection}.

\subsection{Device Temperature Measurement Validation}
\label{sec:ch4-temperature-validation}

The TTC05103JSY 10k$\Omega$ thermistor was compared with a reference infrared
thermometer at the sensor measurement point.
\begin{table}[H]
  \centering
  \caption{Validation results for the device temperature measurement thermistor channel}
  \label{tab:ch4-temperature-validation}
  \begin{tabular}{ccccc}
    \toprule
    \textbf{Test} & \textbf{Device ($^\circ$C)} & \textbf{IR thermometer ($^\circ$C)} & \textbf{Abs. error ($^\circ$C)} & \textbf{Squared error ($^\circ$C$^2$)} \\
    \midrule
    1 & 27.9 & 27.4 & 0.5 & 0.25 \\
    2 & 31.1 & 31.8 & 0.7 & 0.49 \\
    3 & 37.7 & 36.6 & 1.1 & 1.21 \\
    4 & 41.1 & 41.9 & 0.8 & 0.64 \\
    5 & 48.2 & 47.3 & 0.9 & 0.81 \\
    6 & 51.2 & 52.5 & 1.3 & 1.69 \\
    7 & 59.0 & 58.1 & 0.9 & 0.81 \\
    8 & 61.6 & 63.2 & 1.6 & 2.56 \\
    9 & 67.8 & 66.8 & 1.0 & 1.00 \\
    10 & 68.3 & 69.7 & 1.4 & 1.96 \\
    \bottomrule
  \end{tabular}
\end{table}

Table~\ref{tab:ch4-temperature-validation} supports the use of the thermistor channel as a
stable observable input for the thermal estimator. The device temperature measurement channel produced an MAE
of 1.02°C, MSE of 1.142°C$^2$, and RMSE of 1.07°C. This
validates the thermistor channel itself, but it must not be interpreted as a
direct validation of internal contact temperature. 
Hence, it reinforces the
methodological need for the compensated socket-temperature model because the
sensor is external to the highest-temperature contact interface.
The heating tests later
show that the socket terminal can be much hotter than the external thermistor reading
or compensated estimate. 
